\documentclass[letterpaper, 10 pt, conference]{ieeeconf}
\IEEEoverridecommandlockouts
\overrideIEEEmargins

\usepackage{amsmath,amssymb,amsfonts}
\usepackage{graphicx,booktabs}
\usepackage{cite}

\title{\bf Adversarial Commit-Time Governance: Robustness Under Attack, Deception, and Distribution Shift}
\author{[Authors anonymized for submission]}

\begin{document}
\maketitle
\thispagestyle{empty}
\pagestyle{empty}

%%%%%%%%%%%%%%%%%%%%%%%%%%%%%%%%%%%%%%%%%%%%%%%%%%%%%%%%%%%%%%%%%%%%%%%%%%%%%%%%
\begin{abstract}
We analyze the robustness of commit-time governance under adversarial conditions. Four attack classes are considered: state deception, lag inflation, boundary surfing, and recovery disruption. Under bounded perturbations within a calibration envelope, the commit gate preserves soundness with respect to a certified safe subset. Outside this envelope, the system degrades in a fail-closed manner, resulting in denial or safe-mode behavior rather than unsafe execution. Residual risks are characterized as false negatives and reduced performance rather than silent catastrophic failure.
\end{abstract}

%%%%%%%%%%%%%%%%%%%%%%%%%%%%%%%%%%%%%%%%%%%%%%%%%%%%%%%%%%%%%%%%%%%%%%%%%%%%%%%%
\section{Introduction}

Commit-time governance enforces admissibility at the point of action. Prior work established that the gate is sound with respect to a certified subset of recoverable states. However, real systems operate under adversarial conditions, including corrupted observations, delayed signals, and targeted disruptions.

This paper evaluates how the commit gate behaves under such conditions. The focus is not absolute security, but whether failures preserve safety.

%%%%%%%%%%%%%%%%%%%%%%%%%%%%%%%%%%%%%%%%%%%%%%%%%%%%%%%%%%%%%%%%%%%%%%%%%%%%%%%%
\section{Threat Model}

\begin{definition}[Adversarial Perturbation]
The adversary may perturb:
\begin{align}
\hat{\mathbf{x}} &= \mathbf{x} + \delta_x, \quad \|\delta_x\| \le \Delta_x \\
\hat{\tau} &= \tau + \delta_\tau
\end{align}
\end{definition}

\begin{definition}[Calibration Envelope]
\begin{equation}
\mathcal{E}_{cal} =
\left\{
(\delta_x, \delta_\tau) :
\|\delta_x\| \le \Delta_x^{cal},\;
0 \le \delta_\tau \le \Delta_\tau^{cal}
\right\}
\end{equation}
\end{definition}

Within this envelope, barrier parameters remain valid.

%%%%%%%%%%%%%%%%%%%%%%%%%%%%%%%%%%%%%%%%%%%%%%%%%%%%%%%%%%%%%%%%%%%%%%%%%%%%%%%%
\section{Perturbed Commit Evaluation}

The gate evaluates:
\begin{equation}
\text{ALLOW}(u;\hat{\mathbf{x}},\hat{\tau})
\end{equation}

rather than the true state.

We analyze the resulting deviation in admissibility.

%%%%%%%%%%%%%%%%%%%%%%%%%%%%%%%%%%%%%%%%%%%%%%%%%%%%%%%%%%%%%%%%%%%%%%%%%%%%%%%%
\section{Attack Classes}

\subsection{State Deception}

Perturbation:
\begin{equation}
\hat{\mathbf{x}} = \mathbf{x} + \delta_x
\end{equation}

If $\|\delta_x\| \le \Delta_x^{cal}$, then:
\begin{equation}
\mathbf{x} \in \mathcal{S}_{cert}(\tau)
\Rightarrow
\hat{\mathbf{x}} \in \mathcal{S}_{cert}^+(\tau)
\end{equation}

where $\mathcal{S}_{cert}^+$ is an expanded certified subset.

\subsection{Lag Inflation}

\begin{equation}
\hat{\tau} = \tau + \delta_\tau
\end{equation}

Since recoverability degrades monotonically with $\tau$, inflated lag yields stricter conditions:
\begin{equation}
\mathcal{S}_{cert}(\hat{\tau}) \subseteq \mathcal{S}_{cert}(\tau)
\end{equation}

\subsection{Boundary Surfing}

Adversary applies a sequence of small admissible actions:
\begin{equation}
\mathbf{x}_{k+1} = \Phi(\mathbf{x}_k, u_k)
\end{equation}

such that $\mathbf{x}_k \to \partial \mathcal{S}_{cert}$.

This does not violate pointwise admissibility but may approach instability.

\subsection{Recovery Disruption}

During recovery:
\begin{equation}
\dot{\mathbf{x}} = f(\mathbf{x}) + u_{rec} + d_{rec}
\end{equation}

If disturbance satisfies:
\begin{equation}
\|d_{rec}\| < \lambda_{min}(K_{rec}) \|\mathbf{e}\|
\end{equation}

then contraction toward equilibrium is preserved.

%%%%%%%%%%%%%%%%%%%%%%%%%%%%%%%%%%%%%%%%%%%%%%%%%%%%%%%%%%%%%%%%%%%%%%%%%%%%%%%%
\section{Robustness Results}

\begin{theorem}[Adversarial Soundness]
If $(\delta_x,\delta_\tau) \in \mathcal{E}_{cal}$, then:
\begin{equation}
\text{ALLOW}(u;\hat{\mathbf{x}},\hat{\tau})
\Rightarrow
\Phi(\mathbf{x},u) \in \mathcal{S}_{cert}^+(\tau)
\subseteq
\mathcal{R}_{rob}(\tau)
\end{equation}
\end{theorem}

\textit{Interpretation.}
Within the calibration envelope, perturbations preserve soundness up to a bounded expansion.

\begin{proposition}[Fail-Closed Behavior]
If perturbations exceed $\mathcal{E}_{cal}$:
\begin{itemize}
\item State deception: may bypass gate if undetected
\item Lag inflation: leads to conservative denial
\item Recovery disruption: triggers repeated denial or safe mode
\end{itemize}
\end{proposition}

%%%%%%%%%%%%%%%%%%%%%%%%%%%%%%%%%%%%%%%%%%%%%%%%%%%%%%%%%%%%%%%%%%%%%%%%%%%%%%%%
\section{Residual Risk}

Residual risks include:
\begin{itemize}
\item False negatives (safe actions denied)
\item Performance degradation
\item Safe-mode activation
\end{itemize}

Notably, unsafe actions are not silently admitted within the calibration envelope.

%%%%%%%%%%%%%%%%%%%%%%%%%%%%%%%%%%%%%%%%%%%%%%%%%%%%%%%%%%%%%%%%%%%%%%%%%%%%%%%%
\section{Discussion}

The commit gate is not adversarially complete. It guarantees pointwise safety but does not enforce trajectory-level constraints. Boundary surfing remains a limitation without additional monitoring.

The key property is fail-closed degradation: failures result in denial rather than unsafe execution.

%%%%%%%%%%%%%%%%%%%%%%%%%%%%%%%%%%%%%%%%%%%%%%%%%%%%%%%%%%%%%%%%%%%%%%%%%%%%%%%%
\section{Conclusion}

We showed that commit-time governance remains sound under bounded adversarial perturbations and fails safely outside that regime. This establishes robustness as a property of controlled degradation rather than absolute resistance.

\end{document}
